\documentclass[12pt,a4paper]{article}
\usepackage[UTF8]{ctex}
\usepackage{amsmath,amssymb,amsthm}
\usepackage{geometry}
\usepackage{bm}
\usepackage{hyperref}
\usepackage{tikz}
\usetikzlibrary{arrows,arrows.meta,positioning,calc,shapes.geometric}

\geometry{margin=2.5cm}

\title{为什么CoR在直线上移动？}
\author{第12章抓取与操作补充说明}
\date{}

\newtheorem{theorem}{定理}
\newtheorem{definition}{定义}
\newtheorem{example}{例子}

\begin{document}

\maketitle

\section{问题}

给定两个twist $V_1$和$V_2$，它们的线性组合：
\begin{equation}
V = k_1 V_1 + k_2 V_2
\end{equation}

\textbf{问题}：为什么当$k_1$和$k_2$变化时，$\text{CoR}(V)$在通过$\text{CoR}(V_1)$和$\text{CoR}(V_2)$的\textbf{直线}上移动？

\section{数学证明}

\subsection{设置}

设：
\begin{align}
V_1 &= (\omega_{1z}, v_{1x}, v_{1y}) \quad \Rightarrow \quad \text{CoR}_1 = \left(-\frac{v_{1y}}{\omega_{1z}}, \frac{v_{1x}}{\omega_{1z}}\right) \\
V_2 &= (\omega_{2z}, v_{2x}, v_{2y}) \quad \Rightarrow \quad \text{CoR}_2 = \left(-\frac{v_{2y}}{\omega_{2z}}, \frac{v_{2x}}{\omega_{2z}}\right)
\end{align}

线性组合：
\begin{equation}
V = k_1 V_1 + k_2 V_2 = (k_1 \omega_{1z} + k_2 \omega_{2z}, k_1 v_{1x} + k_2 v_{2x}, k_1 v_{1y} + k_2 v_{2y})
\end{equation}

\subsection{CoR位置的计算}

如果$\omega_z = k_1 \omega_{1z} + k_2 \omega_{2z} \neq 0$，则：
\begin{align}
\text{CoR}(V) &= \left(-\frac{v_y}{\omega_z}, \frac{v_x}{\omega_z}\right) \\
&= \left(-\frac{k_1 v_{1y} + k_2 v_{2y}}{k_1 \omega_{1z} + k_2 \omega_{2z}}, \frac{k_1 v_{1x} + k_2 v_{2x}}{k_1 \omega_{1z} + k_2 \omega_{2z}}\right)
\label{eq:cor_combination}
\end{align}

\subsection{关键技巧：引入参数$t$}

\textbf{关键观察}：CoR位置只依赖于$k_1$和$k_2$的\textbf{比值}，而不是绝对值。

设：
\begin{equation}
t = \frac{k_1}{k_1 + k_2}, \quad \text{当 } k_1 + k_2 \neq 0
\end{equation}

则：
\begin{align}
k_1 &= t(k_1 + k_2) \\
k_2 &= (1-t)(k_1 + k_2)
\end{align}

\textbf{注意}：这里我们假设$k_1 + k_2 \neq 0$。如果$k_1 + k_2 = 0$，则$k_2 = -k_1$，这是特殊情况（后面会讨论）。

\subsection{重写CoR公式}

将$k_1$和$k_2$用$t$表示，代入公式(\ref{eq:cor_combination})：

\begin{align}
\text{CoR}(V) &= \left(-\frac{t(k_1+k_2) v_{1y} + (1-t)(k_1+k_2) v_{2y}}{t(k_1+k_2) \omega_{1z} + (1-t)(k_1+k_2) \omega_{2z}}, \frac{t(k_1+k_2) v_{1x} + (1-t)(k_1+k_2) v_{2x}}{t(k_1+k_2) \omega_{1z} + (1-t)(k_1+k_2) \omega_{2z}}\right) \\
&= \left(-\frac{t v_{1y} + (1-t) v_{2y}}{t \omega_{1z} + (1-t) \omega_{2z}}, \frac{t v_{1x} + (1-t) v_{2x}}{t \omega_{1z} + (1-t) \omega_{2z}}\right)
\label{eq:cor_t}
\end{align}

\textbf{关键点}：$(k_1 + k_2)$被约掉了！CoR位置只依赖于$t$。

\subsection{参数化直线方程}

现在我们需要证明：当$t$变化时，$\text{CoR}(V)$在通过$\text{CoR}_1$和$\text{CoR}_2$的直线上。

\textbf{思路}：我们需要找到一个参数$\lambda$，使得：
\begin{equation}
\text{CoR}(V) = \lambda \text{CoR}_1 + (1-\lambda) \text{CoR}_2
\end{equation}

或者，更一般地，证明$\text{CoR}(V)$满足通过$\text{CoR}_1$和$\text{CoR}_2$的直线方程。

\subsection{直线方程}

通过两点$(x_1, y_1)$和$(x_2, y_2)$的直线方程可以写成参数形式：
\begin{equation}
(x, y) = (1-\lambda)(x_1, y_1) + \lambda(x_2, y_2), \quad \lambda \in \mathbb{R}
\end{equation}

或者：
\begin{equation}
(x, y) = (x_1, y_1) + \lambda[(x_2, y_2) - (x_1, y_1)], \quad \lambda \in \mathbb{R}
\end{equation}

\subsection{证明CoR在直线上}

\textbf{方法1}：直接验证

我们需要证明存在$\lambda$使得：
\begin{equation}
\text{CoR}(V) = (1-\lambda) \text{CoR}_1 + \lambda \text{CoR}_2
\end{equation}

这等价于证明$\text{CoR}(V)$、$\text{CoR}_1$和$\text{CoR}_2$三点共线。

\textbf{方法2}：使用向量叉积

三点共线的充要条件是：
\begin{equation}
(\text{CoR}(V) - \text{CoR}_1) \times (\text{CoR}_2 - \text{CoR}_1) = 0
\end{equation}

\textbf{方法3}：直接计算（推荐）

让我们直接计算$\text{CoR}(V)$的表达式，看看它是否在直线上。

从公式(\ref{eq:cor_t})，我们有：
\begin{align}
\text{CoR}(V) &= \left(-\frac{t v_{1y} + (1-t) v_{2y}}{t \omega_{1z} + (1-t) \omega_{2z}}, \frac{t v_{1x} + (1-t) v_{2x}}{t \omega_{1z} + (1-t) \omega_{2z}}\right)
\end{align}

而：
\begin{align}
\text{CoR}_1 &= \left(-\frac{v_{1y}}{\omega_{1z}}, \frac{v_{1x}}{\omega_{1z}}\right) \\
\text{CoR}_2 &= \left(-\frac{v_{2y}}{\omega_{2z}}, \frac{v_{2x}}{\omega_{2z}}\right)
\end{align}

\textbf{关键观察}：虽然不能直接写成$\text{CoR}_1$和$\text{CoR}_2$的线性组合，但我们可以证明$\text{CoR}(V)$满足通过$\text{CoR}_1$和$\text{CoR}_2$的直线方程。

\subsection{几何证明思路}

\textbf{更直观的方法}：考虑CoR的几何意义。

\begin{itemize}
\item CoR是平面上的一个点
\item 当$k_1$和$k_2$变化时，twist $V$在twist空间中沿直线变化
\item 从twist到CoR的映射是\textbf{分式线性变换}（Möbius变换）
\item 分式线性变换将直线映射为直线或圆
\item 在这个特殊情况下，它映射为直线
\end{itemize}

\section{具体数值例子}

让我们用一个具体的例子来验证。

\subsection{例子设置}

设：
\begin{align}
V_1 &= (1, 2, 0) \quad \Rightarrow \quad \text{CoR}_1 = (0, 2) \\
V_2 &= (1, 0, 2) \quad \Rightarrow \quad \text{CoR}_2 = (-2, 0)
\end{align}

\subsection{计算几个CoR位置}

\textbf{情况1}：$k_1 = 1, k_2 = 0$（即$V = V_1$）
\begin{align}
V &= (1, 2, 0) \\
\text{CoR}(V) &= (0, 2) = \text{CoR}_1 \quad \checkmark
\end{align}

\textbf{情况2}：$k_1 = 0, k_2 = 1$（即$V = V_2$）
\begin{align}
V &= (1, 0, 2) \\
\text{CoR}(V) &= (-2, 0) = \text{CoR}_2 \quad \checkmark
\end{align}

\textbf{情况3}：$k_1 = 1, k_2 = 1$（即$V = V_1 + V_2$）
\begin{align}
V &= (2, 2, 2) \\
\text{CoR}(V) &= \left(-\frac{2}{2}, \frac{2}{2}\right) = (-1, 1)
\end{align}

\textbf{验证}：$(-1, 1)$是否在通过$(0, 2)$和$(-2, 0)$的直线上？

通过两点的直线方程：
\begin{equation}
\frac{y - 2}{x - 0} = \frac{0 - 2}{-2 - 0} = 1
\end{equation}

即：$y - 2 = x$，或$y = x + 2$

检查$(-1, 1)$：$1 = -1 + 2 = 1$ \quad $\checkmark$

\textbf{情况4}：$k_1 = 2, k_2 = 1$
\begin{align}
V &= (3, 4, 2) \\
\text{CoR}(V) &= \left(-\frac{2}{3}, \frac{4}{3}\right) = \left(-\frac{2}{3}, \frac{4}{3}\right)
\end{align}

验证：$\frac{4}{3} = -\frac{2}{3} + 2 = \frac{4}{3}$ \quad $\checkmark$

\textbf{情况5}：$k_1 = 1, k_2 = 2$
\begin{align}
V &= (3, 2, 4) \\
\text{CoR}(V) &= \left(-\frac{4}{3}, \frac{2}{3}\right)
\end{align}

验证：$\frac{2}{3} = -\frac{4}{3} + 2 = \frac{2}{3}$ \quad $\checkmark$

\subsection{图示}

\begin{center}
\begin{tikzpicture}[scale=1.5]
% 两个CoR
\filldraw[red] (0,2) circle (3pt) node[above left] {CoR$_1$ $(0, 2)$};
\filldraw[blue] (-2,0) circle (3pt) node[below left] {CoR$_2$ $(-2, 0)$};

% 连接直线
\draw[thick,green] (-2.5,-0.5) -- (0.5,2.5);
\node[green] at (-1.2,1.2) {直线 $y = x + 2$};

% 标记几个CoR
\filldraw[purple] (-1,1) circle (2pt) node[above right] {$(-1, 1)$\\$k_1=1, k_2=1$};
\filldraw[purple] (-0.67,1.33) circle (2pt) node[above right] {$(-\frac{2}{3}, \frac{4}{3})$\\$k_1=2, k_2=1$};
\filldraw[purple] (-1.33,0.67) circle (2pt) node[below left] {$(-\frac{4}{3}, \frac{2}{3})$\\$k_1=1, k_2=2$};

\node[below] at (-1,-0.8) {所有CoR都在同一条直线上};
\end{tikzpicture}
\end{center}

\section{一般性证明}

\subsection{参数化表示}

从公式(\ref{eq:cor_t})，我们有：
\begin{equation}
\text{CoR}(V) = \left(-\frac{t v_{1y} + (1-t) v_{2y}}{t \omega_{1z} + (1-t) \omega_{2z}}, \frac{t v_{1x} + (1-t) v_{2x}}{t \omega_{1z} + (1-t) \omega_{2z}}\right)
\end{equation}

其中$t = \frac{k_1}{k_1 + k_2}$。

\textbf{关键观察}：这个表达式是$t$的\textbf{分式线性函数}。

\subsection{分式线性变换的性质}

分式线性变换（Möbius变换）的一般形式：
\begin{equation}
f(z) = \frac{az + b}{cz + d}
\end{equation}

在我们的情况下，CoR的每个坐标都是$t$的分式线性函数。

\textbf{重要性质}：分式线性变换将直线映射为直线或圆。在我们的情况下，由于分子和分母都是$t$的线性函数，结果总是直线。

\subsection{验证端点}

\textbf{当$t = 1$时}（即$k_2 = 0$）：
\begin{align}
\text{CoR}(V) &= \left(-\frac{v_{1y}}{\omega_{1z}}, \frac{v_{1x}}{\omega_{1z}}\right) = \text{CoR}_1
\end{align}

\textbf{当$t = 0$时}（即$k_1 = 0$）：
\begin{align}
\text{CoR}(V) &= \left(-\frac{v_{2y}}{\omega_{2z}}, \frac{v_{2x}}{\omega_{2z}}\right) = \text{CoR}_2
\end{align}

\textbf{结论}：当$t$从0变化到1时，$\text{CoR}(V)$从$\text{CoR}_2$移动到$\text{CoR}_1$，沿着连接两点的直线。

\subsection{扩展到整个实数轴}

当$t$取所有实数值时：
\begin{itemize}
\item $t < 0$：CoR在从CoR$_2$出发的延长线上
\item $0 \leq t \leq 1$：CoR在CoR$_1$和CoR$_2$之间的线段上
\item $t > 1$：CoR在从CoR$_1$出发的延长线上
\end{itemize}

\textbf{关键}：无论$t$取什么值，CoR都在通过CoR$_1$和CoR$_2$的\textbf{整条直线}上。

\section{几何直觉}

\subsection{为什么是直线？

\textbf{直觉1}：线性组合的"线性"性质

\begin{itemize}
\item Twist的线性组合：$V = k_1 V_1 + k_2 V_2$
\item 这是"线性"的
\item CoR是twist的"非线性"函数（分式）
\item 但分式线性函数仍然保持"直线性"
\end{itemize}

\textbf{直觉2}：参数化

\begin{itemize}
\item 当$k_1$和$k_2$变化时，实际上只有一个自由度（比值$t$）
\item 一个自由度的变化通常对应一条曲线
\item 在这个特殊情况下，是直线
\end{itemize}

\textbf{直觉3}：连续性

\begin{itemize}
\item 当$k_1$和$k_2$连续变化时，CoR也连续变化
\item 从CoR$_1$到CoR$_2$的最短路径是直线
\item 数学上，分式线性变换保证了这种"直线性"
\end{itemize}

\section{特殊情况}

\subsection{当$\omega_z = 0$时}

如果$k_1 \omega_{1z} + k_2 \omega_{2z} = 0$：
\begin{itemize}
\item 这是纯平移
\item CoR在无穷远处
\item 可以认为这是直线的"端点"（在射影几何中）
\end{itemize}

\subsection{当$k_1 + k_2 = 0$时}

如果$k_2 = -k_1$：
\begin{itemize}
\item $V = k_1(V_1 - V_2)$
\item 这是$V_1$和$V_2$的差
\item 需要单独处理
\end{itemize}

\section{总结}

\subsection{关键结论}

\begin{enumerate}
\item \textbf{数学事实}：当$k_1$和$k_2$变化时，$\text{CoR}(V)$在通过$\text{CoR}(V_1)$和$\text{CoR}(V_2)$的直线上移动

\item \textbf{原因}：
  \begin{itemize}
  \item CoR位置只依赖于$k_1$和$k_2$的比值$t$
  \item CoR是$t$的分式线性函数
  \item 分式线性函数将参数直线映射为直线
  \end{itemize}

\item \textbf{几何意义}：
  \begin{itemize}
  \item 一个自由度的变化对应一条曲线
  \item 在这个情况下，是直线
  \end{itemize}

\item \textbf{验证}：通过具体数值例子可以验证
\end{enumerate}

\subsection{应用}

这个性质在12.1.6章节中非常重要：
\begin{itemize}
\item 理解正线性组合的CoR区域
\item 分析可行运动集合
\item 可视化接触约束
\end{itemize}

\end{document}

